\section{Baza podataka}
Sembot trenutno koristi JSON datoteku (qa\_data.json) kao jednostavnu bazu podataka za pohranu pitanja i odgovora. Struktura datoteke je sljedeća:

\begin{minted}[fontsize=\small, linenos, breaklines, frame=lines, bgcolor=gray!10]{json}
{
  "qa_pairs": [
    {
      "question": "Koji su rokovi za prijavu ispita?",
      "answer": "Rokovi za prijavu ispita su najkasnije 3 dana prije termina ispita putem studomata."
    },
    {
      "question": "Kada počinje akademska godina?",
      "answer": "Za detaljni pregled svih važnih datuma i događaja akademske godine, pogledajte akademski kalendar na: https://www.unipu.hr/akademski_kalendar"
    },
    {
      "question": "Koliko iznose ECTS bodovi za završetak studija?",
      "answer": "Za završetak preddiplomskog studija potrebno je ostvariti 180 ECTS bodova, a za diplomski studij dodatnih 120 ECTS bodova."
    },
    {
      "question": "Kako mogu kontaktirati studentsku službu?",
      "answer": "Studentsku službu možete kontaktirati putem emaila ured-za-studente@unipu.hr ili telefonom na broj +385 52 377 000 radnim danom od 10-14h."
    },
    {
      "question": "Koji su uvjeti za upis više godine?",
      "answer": "Za upis više godine studija potrebno je ostvariti minimalno 50 ECTS bodova iz tekuće akademske godine."
    },
    ...
  ]
}
\end{minted}

\subsection{Prednosti i nedostaci JSON baze}

\paragraph{Prednosti JSON baze}
\begin{itemize}
    \item \textbf{Jednostavnost:} JSON datoteke su lako čitljive i uređive, što olakšava dodavanje novih pitanja i odgovora.
    \item \textbf{Brzina:} Mala veličina datoteke omogućuje brzo učitavanje i pretraživanje.
    \item \textbf{Nema potrebe za DBMS:} Ne zahtijeva instalaciju složenih sustava upravljanja bazama podataka.
    \item \textbf{Portabilnost:} JSON datoteke mogu se lako prenijeti na različite sustave.
\end{itemize}

\paragraph{Nedostaci JSON baze}
\begin{itemize}
    \item \textbf{Skalabilnost:} Za velike količine podataka (npr. tisuće pitanja), JSON može postati neefikasan zbog potrebe za učitavanjem cijele datoteke u memoriju.
    \item \textbf{Ograničene mogućnosti upita:} JSON ne podržava složene upite poput SQL-a, što ograničava funkcionalnost za napredne analize.
    \item \textbf{Sigurnost:} Bez dodatnih mjera, JSON datoteke su ranjive na neovlaštene promjene.
\end{itemize}

\subsection{Upravljanje podacima}
Trenutno se podaci u qa\_data.json ručno uređuju, ali u budućnosti može se razviti administrativno sučelje za dodavanje i uređivanje pitanja i odgovora. Ovo sučelje moglo bi biti implementirano kao dio web aplikacije, omogućujući administratorima jednostavno upravljanje sadržajem.